% eksperymentalne tłumaczenie na włoski

%

\index{punkt!Lemoine'a} % lang-pl
\index{punto!di Lemoine} % lang-it

Emile Lemoine odczyta w 1873 roku pracę \emph{,,Sur quelques propriétés d'un point remarquable du triangle''}, gdzie po raz pierwszy pojawią się bez dowodów fakty \ref{random_prp_symmedians_proportional_squares} i \ref{punkt_lemoine}. % lang-pl
Termin symediana zaproponuje Maurice d'Ocagne w 1883 roku, jako zastępstwo \emph{antyrównoległej środkowej} używanego przez Lemoine'a. % lang-pl
Nel 1873 Émile Lemoine presenterà il lavoro \emph{,,Sur quelques propriétés d'un point remarquable du triangle''}, nel quale compariranno per la prima volta, senza dimostrazione, i fatti \ref{random_prp_symmedians_proportional_squares} e \ref{punkt_lemoine}. % lang-it
Il termine \emph{symmediana} sarà proposto da Maurice d'Ocagne nel 1883 come sostituto di \emph{mediana antiparallela}, espressione usata da Lemoine. % lang-it
% Nouvelles Annales de Mathématiques, 3rd series, II, 451 (1883).
Punkt Lemoine'a pojawia się u Józefa Neuberga (praca \emph{Mémoire sur le tétraèdre} z 1886 roku). % lang-pl
Il punto di Lemoine compare nell'opera di Joseph Neuberg (\emph{Mémoire sur le tétraèdre}, 1886). % lang-it
% https://www.persee.fr/doc/marb_0770-8459_1886_num_37_1_2317
W~1847 roku Ernst Grebe skonstruuje kwadraty na zewnątrz trójkąta i odkryje punkt Lemoine'a przed Lemoinem, dlatego w literaturze padać będzie też nazwa punkt Grebego. % lang-pl
Nel 1847 Ernst Grebe costruirà quadrati all'esterno di un triangolo e scoprirà il punto di Lemoine prima dello stesso Lemoine; per questo motivo nella letteratura compare anche il nome di punto di Grebe. % lang-it

Wszystko to pięknie opisze Altshiller-Court \cite{altshiller_2007}. % lang-pl
Tutto ciò è descritto magnificamente da Altshiller-Court \cite{altshiller_2007}. % lang-it

% UW: Twierdzenie Cevy (wraz z trygonometryczną wersją),
% przykłady punktów szczególnych trójkąta:
% punkt Nagela (Guzicki-4),
% punkt Gergonne'a (guzicki4),

% https://www.deltami.edu.pl/2016/05/izogonalne-sprzezenie-i-symediany/

\begin{definition}
[symediana] % lang-pl
[symmediana] % lang-it
    Odbicie symetryczne środkowej trójkąta względem dwusiecznej wychodzącej z~tego samego wierzchołka nazywamy symedianą. % lang-pl
    La riflessione della mediana di un triangolo rispetto alla bisettrice uscente dallo stesso vertice si chiama symmediana. % lang-it
\end{definition}

Symediana to prosta izogonalna ze środkową (Zetel \cite[s. 111]{zetel_2020}; Ambroszczyk i Komisarczyk \cite[s. 26]{komisarczyk_2024}; Dutka w $\Delta_{15}^2$). % lang-pl
La symmediana è la retta izogonale della mediana (Zetel \cite[p. 111]{zetel_2020}; Ambroszczyk e Komisarczyk \cite[p. 26]{komisarczyk_2024}; Dutka in $\Delta_{15}^2$). % lang-it
% https://www.deltami.edu.pl/2015/02/krotka-opowiesc-o-symedianie/
Symediana jest miejscem geometrycznym punktów, których odległości od dwóch boków trójkąta są proporcjonalne do tych boków. % lang-pl
Ale dużo ważniejszą jej własnością jest: % lang-pl
La symmediana è il luogo geometrico dei punti le cui distanze da due lati del triangolo sono proporzionali a tali lati. % lang-it
Ma una proprietà molto più importante è la seguente: % lang-it

\begin{proposition}
    \label{random_prp_symmedians_proportional_squares}
    Symediana poprowadzona z wierzchołka $C$ trójkąta $ABC$ dzieli wewnętrznie przeciwległy bok proporcjonalnie do kwadratów długości boków przyległych: % lang-pl
    La symmediana condotta dal vertice $C$ del triangolo $ABC$ divide internamente il lato opposto in proporzione ai quadrati delle lunghezze dei lati adiacenti: % lang-it
    \begin{equation}
        \frac{AF}{BF} = \frac{b^2}{a^2}.
    \end{equation}
\end{proposition}

I odwrotnie, punkt leżący na boku $AB$ i spełniający powyższą proporcję wyznacza symedianę z~wierzchołka $C$. % lang-pl
Zetel \cite[s. 111, 112]{zetel_2020} uważa to za szczególny przypadek twierdzenie Steinera o~prostych izogonalnych. % lang-pl
(Chyba nie ma tego twierdzenia w tej książce). % lang-pl
Viceversa, un punto appartenente al lato $AB$ che soddisfa la proporzione precedente determina la symmediana uscente dal vertice $C$. % lang-it
Zetel \cite[p. 111, 112]{zetel_2020} considera questo un caso particolare del teorema di Steiner sulle rette izogonali. % lang-it
(Probabilmente questo teorema non compare nel presente libro). % lang-it

\begin{proposition}
    Symediana wychodząca z wierzchołka $A$ trójkąta $ABC$ ma długość % lang-pl
    La symmediana uscente dal vertice $A$ del triangolo $ABC$ ha lunghezza % lang-it
    \begin{equation}
        \frac{bc}{b^2 + c^2} \cdot \sqrt {2b^2 + 2c^2 - a^2}.
    \end{equation}
\end{proposition}

Zetel \cite[s. 119]{zetel_2020} pisze, że można dostać tę długość z twierdzenia Stewarta. % lang-pl
Zetel \cite[p. 119]{zetel_2020} scrive che questa lunghezza può essere ottenuta mediante il teorema di Stewart. % lang-it
% ale nie trzeba.

\begin{proposition}
[punkt Lemoine'a] % lang-pl
[punto di Lemoine] % lang-it
    \label{punkt_lemoine}%
    W ustalonym trójkącie $ABC$, trzy symediany przecinają się w jednym punkcie $L$, zwanym punktem Lemoine'a. % lang-pl
    In un triangolo fissato $ABC$, le tre symmediane si incontrano in un unico punto $L$, detto punto di Lemoine. % lang-it
\end{proposition}

Zetel \cite[s. 114]{zetel_2020} sugeruje udowodnić \ref{punkt_lemoine} na podstawie twierdzenia odwrotnego do twierdzenia Cevy. % lang-pl
Ross Honsberger nazwie istnienie punktu Lemoine'a ,,jednym z klejnotów koronnych współczesnej geometrii''. % lang-pl
Zetel \cite[p. 114]{zetel_2020} suggerisce di dimostrare \ref{punkt_lemoine} mediante il teorema inverso di Ceva. % lang-it
Ross Honsberger definirà l'esistenza del punto di Lemoine come \emph{uno dei gioielli della corona della geometria moderna}. % lang-it

\begin{proposition}
    Punkt Lemoine'a minimalizuje sumę kwadratów odległości od boków trójkąta. % lang-pl
    Il punto di Lemoine minimizza la somma dei quadrati delle distanze dai lati del triangolo. % lang-it
\end{proposition}

% TODO: https://www.deltami.edu.pl/media/articles/2015/02/delta-2015-02-krotka-opowiesc-o-symedianie.pdf

\begin{proposition}
    Punkt Lemoine'a $L$ trójkąta $ABC$ jest środkiem ciężkości trójkąta, którego wierzchołki to rzuty punktu $L$ na boki $AB$, $BC$, $AC$. % lang-pl
    Il punto di Lemoine $L$ del triangolo $ABC$ è il baricentro del triangolo i cui vertici sono le proiezioni di $L$ sui lati $AB$, $BC$, $AC$. % lang-it
\end{proposition}

\begin{proposition}
[twierdzenie Schlömilcha] % lang-pl
[teorema di Schlömilch] % lang-it
    Trzy proste łączące środki boków trójkąta ze środkami odpowiednich wysokości są współpękowe, przechodzą przez punkt Lemoine'a. % lang-pl
    Le tre rette che congiungono i punti medi dei lati di un triangolo con i punti medi delle rispettive altezze sono concorrenti e passano per il punto di Lemoine. % lang-it
\end{proposition}

\index{twierdzenie!Schlömilcha} % lang-pl
\index{teorema!di Schlömilch} % lang-it

Napisze o tym duet Bogdańska, Neugebauer \cite[s. 195]{neugebauer_2018} (jako ćwiczenie, bez wzmianki o nazwie punktu), ale też Zetel \cite[s. 34]{zetel_2020}. % lang-pl
Ne scrivono Bogdańska e Neugebauer \cite[p. 195]{neugebauer_2018} (come esercizio, senza menzionare il nome del punto), ma anche Zetel \cite[p. 34]{zetel_2020}. % lang-it
% https://ems.press/content/serial-article-files/51635 VVV
Twierdzenie napotka w 1862 roku Oskar Schlömilch, cokolwiek to znaczy. % lang-pl
Il teorema fu scoperto da Oskar Schlömilch nel 1862, qualunque cosa ciò voglia dire. % lang-it

\begin{definition}
[antyrównoległa] % lang-pl
[antiparallela] % lang-it
    Antyrównoległa do boku $AB$ w trójkącie $ABC$ to prosta, która tnie bok naprzeciw wierzchołka $A$ pod takim samym kątem jak ten przy wierzchołku $A$ oraz bok naprzeciw wierzchołka $B$ pod takim samym kątem jak ten przy wierzchołku $B$. % lang-pl
    Un'antiparallela al lato $AB$ nel triangolo $ABC$ è una retta che interseca il lato opposto al vertice $A$ con lo stesso angolo presente in $A$ e il lato opposto al vertice $B$ con lo stesso angolo presente in $B$. % lang-it
\end{definition}

\begin{proposition}
    Trzy antyrównoległe do boków trójkąta przechodzące przez punkt Lemoine'a wyznaczają sześć nowych punktów na brzegu trójkąta. Leżą one na pierwszym okręgu Lemoine'a. % lang-pl
    Le tre antiparallele ai lati del triangolo passanti per il punto di Lemoine determinano sei nuovi punti sul contorno del triangolo. Essi giacciono sul primo cerchio di Lemoine. % lang-it
\end{proposition}

\begin{proposition}
    Trzy równoległe do boków trójkąta przechodzące przez punkt Lemoine'a wyznaczają sześć nowych punktów na brzegu trójkąta. Leżą one na drugim okręgu Lemoine'a. % lang-pl
    Le tre rette parallele ai lati del triangolo passanti per il punto di Lemoine determinano sei nuovi punti sul contorno del triangolo. Essi giacciono sul secondo cerchio di Lemoine. % lang-it
\end{proposition}

Istnieje jeszcze trzeci okrąg Lemoine'a, ale jego konstrukcja nie jest już taka prosta. % lang-pl
Esiste anche un terzo cerchio di Lemoine, ma la sua costruzione non è altrettanto semplice. % lang-it

%